%% mocvrp.tex -- Multi-objective Capacitated Vehicle Routing Problem with
%% Optional Service (MO-CVRP-OS): problem definition and MILP formulation.
%%
%% Builds with
%%     latexmk -pdf mocvrp.tex
%% which runs pdflatex/bibtex as many times as needed.  By hand:
%%     pdflatex mocvrp && bibtex mocvrp && pdflatex mocvrp && pdflatex mocvrp
%% Requires references.bib alongside this file.  No shell-escape, no other
%% generated input files.

\documentclass[11pt,a4paper]{article}

\usepackage[T1]{fontenc}
\usepackage[utf8]{inputenc}
\usepackage{lmodern}
\usepackage{microtype}
\usepackage[margin=2.6cm]{geometry}

\usepackage{amsmath}
\usepackage{amssymb}
\usepackage{amsthm}
\usepackage{mathtools}
\usepackage{booktabs}
\usepackage{enumitem}
\usepackage{array}
\usepackage{siunitx}
\sisetup{group-separator={,},group-minimum-digits=4}

\usepackage[hidelinks]{hyperref}
\usepackage[capitalise,noabbrev]{cleveref}

\allowdisplaybreaks

\theoremstyle{definition}
\newtheorem{definition}{Definition}[section]
\newtheorem{assumption}[definition]{Assumption}
\newtheorem{example}[definition]{Example}
\newtheorem{remark}[definition]{Remark}

\crefname{assumption}{Assumption}{Assumptions}
\Crefname{assumption}{Assumption}{Assumptions}
\crefname{example}{Example}{Examples}
\Crefname{example}{Example}{Examples}
\crefname{remark}{Remark}{Remarks}
\Crefname{remark}{Remark}{Remarks}

\theoremstyle{plain}
\newtheorem{proposition}[definition]{Proposition}
\newtheorem{lemma}[definition]{Lemma}
\newtheorem{corollary}[definition]{Corollary}

% ---- notation shortcuts -------------------------------------------------
\newcommand{\R}{\mathbb{R}}
\newcommand{\Z}{\mathbb{Z}}
\newcommand{\Nat}{\mathbb{N}}
\newcommand{\Feas}{\mathcal{F}}
\newcommand{\Routes}{\mathcal{R}}
\newcommand{\Ground}{\mathcal{S}}
\newcommand{\Arcs}{A}
\newcommand{\Cust}{N}
\newcommand{\Nodes}{V}
\newcommand{\Veh}{K}
\newcommand{\nint}{\operatorname{nint}}
\newcommand{\Lmax}{L^{\max}}
\newcommand{\Wmax}{W^{\max}}
\newcommand{\Omax}{O^{\max}}
\newcommand{\Dmax}{D^{\max}}
\newcommand{\Lub}{\overline{L}}
\newcommand{\Oub}{\overline{O}}
\DeclareMathOperator*{\argmin}{arg\,min}
\DeclareMathOperator*{\argmax}{arg\,max}

\title{\bfseries A Multi-Objective Capacitated Vehicle Routing Problem\\
with Optional Service:\\ Problem Definition and Mixed-Integer Formulation}
\author{}
\date{\today}

\begin{document}
\maketitle

\begin{abstract}
\noindent
This document defines \textup{\textsc{MO-CVRP-OS}}, a multi-objective extension of the
Capacitated Vehicle Routing Problem (CVRP) in which serving a customer is
\emph{optional}, and proposes a mixed-integer linear programming (MILP)
formulation for it.  The problem is defined directly over the CVRPLIB instances
distributed with this repository (the \textsf{X}, \textsf{XL}, \textsf{AGS} and
\textsf{DIMACS} families, \num{222} instances in total), and every modelling
assumption is checked against the actual instance data.  Five objectives are
considered: maximising the number of delivered orders, minimising the number of
routes, minimising the largest route diameter, maximising route balance in a
max--min (egalitarian) sense, and minimising total travelled distance.  The
first four are the objectives of interest; the fifth is included because,
as shown in \cref{prop:degeneracy}, the first four are provably blind to the
uniform inflation of route lengths.  All non-linear elements of the objective
vector -- a maximum, several minima, and three ratios whose denominators are
themselves decision-dependent maxima -- are linearised exactly, and the
correctness of each linearisation is proved.  Modelling choices that the
informal specification leaves open are made explicit, and the alternatives are
stated and compared.
\end{abstract}

\tableofcontents

% =========================================================================
\section{Introduction}
\label{sec:intro}

The CVRP asks for a set of minimum-cost vehicle routes, each starting and
ending at a single depot, that jointly serve every customer exactly once
without violating a common vehicle capacity~\cite{TothVigo2014}.  It is a
single-objective problem: the cost is almost always the total travelled
distance, and the service requirement is hard -- every customer must be
visited.

Both of these features are restrictive in practice.  A dispatcher who cannot
serve all orders on a given day would rather choose \emph{which} orders to
postpone than declare the instance infeasible, and total distance is rarely the
only criterion: fleet size, the geographic compactness of each route, and the
equitable distribution of work across drivers all matter and conflict with one
another.  This document formalises one particular combination of these
concerns.

\paragraph{Contributions and scope.}
\Cref{sec:data} fixes the interpretation of the CVRPLIB instance files shipped
in \texttt{instances/} and records, as standing assumptions, the properties of
that data that the later proofs rely on.  \Cref{sec:problem} defines the
feasible set at the level of solutions, independently of any formulation.
\Cref{sec:objectives} defines the five objectives, together with the
conventions that make them well defined in degenerate cases.
\Cref{sec:pareto} recalls the required multi-objective terminology.
\Cref{sec:milp} gives the MILP formulation and proves it correct.
\Cref{sec:variants} discusses the modelling decisions that were open and the
alternatives to the ones taken.

Multi-objective vehicle routing has a substantial literature of its own; see
\cite{Jozefowiez2008} for a survey of the objectives that have been combined
with routing and of the solution methods used.

This document is a \emph{formulation} document.  It does not propose a solution
algorithm, does not report computational experiments, and does not define a
benchmarking protocol.  The model of \cref{sec:milp} is a reference
formulation: it is exactly solvable only on small instances, whereas the
instances in \texttt{instances/} have up to \num{30000} customers.

% =========================================================================
\section{Instance data and CVRPLIB conventions}
\label{sec:data}

\subsection{The instance files}

The repository contains \num{222} instances in the TSPLIB-inspired CVRPLIB
format, in four families: \textsf{X} (\num{100} instances,
\cite{Uchoa2017}), \textsf{XL} (\num{100} instances, \cite{Queiroga2026}),
\textsf{AGS} (\num{10} instances, \cite{Arnold2019}) and \textsf{DIMACS}
(\num{12} instances contributed to the 12th DIMACS Implementation Challenge,
\cite{DIMACS2021}).  Every file declares \texttt{TYPE~:~CVRP}.

A file with \texttt{DIMENSION~:~$n+1$} describes $n+1$ locations numbered
$1,\dots,n+1$.  The \texttt{DEPOT\_SECTION} of every instance in the repository
contains the single entry $1$, so location $1$ is the depot and locations
$2,\dots,n+1$ are the customers.  Throughout this document the depot is
re-indexed to $0$ and customer $p$ of the file becomes customer $p-1$ of the
model, so that
\begin{equation}
  \Nodes \;=\; \{0\} \cup \Cust, \qquad \Cust \;=\; \{1,2,\dots,n\},
  \qquad n \;=\; \texttt{DIMENSION} - 1 .
  \label{eq:nodes}
\end{equation}
The \texttt{CAPACITY} field gives the common vehicle capacity $Q \in \Z_{>0}$,
and the \texttt{DEMAND\_SECTION} gives the demands $q_i \in \Z_{\ge 0}$, with
$q_0 = 0$ for the depot in all \num{222} instances.

\subsection{Distances}
\label{sec:distances}

Two \texttt{EDGE\_WEIGHT\_TYPE} values occur, and they must be handled
differently.

\paragraph{\texttt{EUC\_2D} (\num{210} instances).}
The \texttt{NODE\_COORD\_SECTION} gives coordinates $(a_i,b_i)$ and the
distance is the Euclidean distance \emph{rounded to the nearest integer},
following the TSPLIB95 convention:
\begin{equation}
  d_{ij} \;=\; \nint\!\Bigl( \sqrt{(a_i-a_j)^2 + (b_i-b_j)^2} \Bigr),
  \qquad \nint(t) = \bigl\lfloor t + \tfrac{1}{2} \bigr\rfloor .
  \label{eq:euc2d}
\end{equation}
The instances of the literature that deviate from this convention by using
unrounded distances -- \textsf{CMT4}, \textsf{CMT5}, \textsf{tai385} and
\textsf{Golden9}--\textsf{Golden20}, as listed in \cite{DIMACS2021} -- are not
part of this repository, so \eqref{eq:euc2d} applies uniformly to all
\num{210} \texttt{EUC\_2D} instances here.

\paragraph{\texttt{EXPLICIT} with \texttt{EDGE\_WEIGHT\_FORMAT~:~LOWER\_ROW}
(\num{12} instances).}
The six \textsf{Loggi} and six \textsf{ORTEC} instances give an explicit
integer distance matrix, listed row by row below the diagonal, i.e.\ row $p$
contains $d_{p,1},\dots,d_{p,p-1}$.  The \texttt{LOWER\_ROW} format encodes a
symmetric matrix.  These are \emph{real road distances}, and they are not
metric: sampling $2\times 10^{5}$ distinct triples per instance finds triangle
inequality violations $d_{ij} > d_{ik} + d_{kj}$ in all twelve, in
$0.51\%$ to $1.66\%$ of the sampled triples for the six \textsf{Loggi}
instances and in $0.011\%$ to $0.054\%$ for the six \textsf{ORTEC} instances.  These instances also carry a
\texttt{NODE\_COORD\_SECTION}, which is informative only: the explicit matrix
takes precedence.

Accordingly, the model below never assumes the triangle inequality, and never
uses a valid inequality that would require it.

\subsection{Fleet size}
\label{sec:fleet}

No instance in the repository declares a fleet size: the only keys that occur
anywhere are \texttt{NAME}, \texttt{COMMENT}, \texttt{TYPE},
\texttt{DIMENSION}, \texttt{EDGE\_WEIGHT\_TYPE},
\texttt{EDGE\_WEIGHT\_FORMAT}, \texttt{NODE\_COORD\_TYPE} and
\texttt{CAPACITY}.  In particular the \texttt{-k} suffix of an instance name is
\emph{descriptive, not prescriptive}: it records the number of routes of a
reference solution at generation time.  This is directly confirmed by the best-known-solution files distributed
alongside the instances: of the \num{112} of them whose name carries a
\texttt{-k} suffix, \num{22} use strictly \emph{more} routes than the name
advertises and none uses fewer -- for instance \texttt{X-n101-k25.sol}
contains \num{26} routes.

The fleet is therefore modelled as homogeneous and unlimited, which is the
standard CVRPLIB reading, and the number of vehicles enters the model as an
\emph{objective} (\cref{sec:objectives}) rather than as a constraint.

\subsection{Standing assumptions}

The following properties were verified by parsing all \num{222} instance files;
they are used in the proofs of \cref{sec:milp} and are collected here so that
the dependency is explicit.

\begin{assumption}[Distance matrix]
\label{as:dist}
$d : \Nodes \times \Nodes \to \Z_{\ge 0}$ satisfies $d_{ii} = 0$ and
$d_{ij} = d_{ji}$.  The triangle inequality is \emph{not} assumed.  Distinct
locations may be at distance zero: this occurs for customer pairs in the ten
\textsf{AGS} instances, which are the only ones containing duplicate
coordinates, and in ten of the twelve \texttt{EXPLICIT} instances, which
contain up to four zero-distance customer pairs each.
\end{assumption}

\begin{assumption}[Demands]
\label{as:demand}
$q_0 = 0$ and $1 \le q_i \le Q$ for every customer $i \in \Cust$.  Over the
\num{222} instances the customer demands range in $[1,100]$ and the capacities
in $[3,\num{10064}]$; no instance contains a customer with zero demand, and no
instance contains a customer whose demand exceeds the vehicle capacity.
\end{assumption}

\begin{assumption}[Depot distances]
\label{as:depot}
$d_{0i} \ge 1$ for every customer $i \in \Cust$.  No customer of any instance
in the repository -- \texttt{EUC\_2D} or \texttt{EXPLICIT} -- lies at distance
zero from the depot.
\end{assumption}

\begin{assumption}[Orders]
\label{as:orders}
Each customer $i \in \Cust$ carries a number of \emph{orders}
$o_i \in \Z_{\ge 1}$.  CVRPLIB associates exactly one demand quantity with each
customer, so the default reading used throughout is $o_i = 1$ for all $i$, in
which case ``number of delivered orders'' and ``number of served customers''
coincide.  The parameter is kept explicit because the formulation is unchanged
for arbitrary positive integer $o_i$ (see \cref{sec:alt-f1}).
\end{assumption}

\Cref{as:demand,as:depot} are not cosmetic.  \Cref{as:demand} is what makes the
single-commodity flow connectivity constraints of \cref{sec:milp} valid
(\cref{prop:validity}), and \cref{as:depot} is what makes the route-balance
objective well defined without an artificial safeguard (\cref{lem:welldef}).
\Cref{as:demand} also guarantees that the trivial single-customer route is
feasible for every customer, so that the ``maximum service'' extreme of the
Pareto front is attained with all customers served.

% =========================================================================
\section{The problem}
\label{sec:problem}

\subsection{Routes and solutions}

Let $G = (\Nodes, \Arcs)$ be the complete digraph on the node set
\eqref{eq:nodes}, with arc set $\Arcs = \{(i,j) \in \Nodes\times\Nodes : i\ne j\}$
and arc lengths $d_{ij}$.

\begin{definition}[Route]
\label{def:route}
A \emph{route} is a sequence $r = (0, i_1, i_2, \dots, i_m, 0)$ with $m \ge 1$
and $i_1,\dots,i_m$ pairwise distinct customers of $\Cust$.  Its
\emph{customer set} is $C(r) = \{i_1,\dots,i_m\}$, and its \emph{length},
\emph{load}, \emph{order count} and \emph{diameter} are
\begin{align}
  L(r) &\;=\; d_{0 i_1} \;+\; \sum_{h=1}^{m-1} d_{i_h i_{h+1}} \;+\; d_{i_m 0},
  \label{eq:Lr}\\
  W(r) &\;=\; \sum_{i \in C(r)} q_i,
  \label{eq:Wr}\\
  O(r) &\;=\; \sum_{i \in C(r)} o_i,
  \label{eq:Or}\\
  D(r) &\;=\; \max_{i,j \in C(r)} d_{ij}
        \;=\; \max \bigl( \{0\} \cup \{ d_{ij} : i,j \in C(r),\, i \ne j \} \bigr).
  \label{eq:Dr}
\end{align}
A route is \emph{capacity feasible} if $W(r) \le Q$.
\end{definition}

Two conventions in \cref{def:route} deserve comment.  First, a route is
required to be non-empty ($m \ge 1$); an ``empty route'' is not a route but the
absence of one, which is what makes the route-counting objective $f_2$ below
meaningful.  Second, the diameter \eqref{eq:Dr} is taken over the customers of
the route \emph{only}: the depot is excluded, so a single-customer route has
diameter $0$.  This is the definition adopted throughout; the depot-inclusive
and depot-radius alternatives are discussed in \cref{sec:alt-diameter}.

\begin{definition}[Solution]
\label{def:solution}
A \emph{solution} is a finite set $\Routes = \{r_1,\dots,r_{|\Routes|}\}$ of
capacity-feasible routes whose customer sets are pairwise disjoint, i.e.\
$C(r) \cap C(r') = \emptyset$ for all $r \ne r' \in \Routes$.  Its
\emph{served set} and \emph{unserved set} are
\begin{equation}
  \Ground(\Routes) \;=\; \bigcup_{r \in \Routes} C(r),
  \qquad
  \Cust \setminus \Ground(\Routes).
\end{equation}
The feasible set of \textup{\textsc{MO-CVRP-OS}} is the set $\Feas$ of all solutions,
including the empty solution $\Routes = \emptyset$.
\end{definition}

Three points distinguish $\Feas$ from the feasible set of the CVRP.

\begin{enumerate}[label=(\roman*),leftmargin=*]
\item \emph{Service is optional.}  \Cref{def:solution} requires the customer
  sets to be disjoint but does not require them to cover $\Cust$.  Customers
  left uncovered are simply not served, and no penalty is charged for them:
  service enters the model through the objective $f_1$ instead.  This places
  \textup{\textsc{MO-CVRP-OS}} in the family of routing problems with profits
  \cite{Archetti2014}, of which the prize-collecting and team-orienteering
  problems are the best known members.
\item \emph{No customer is visited twice.}  Disjointness of the customer sets
  together with the distinctness requirement inside a route means that a served
  customer is visited exactly once, as in the CVRP.  Split delivery is not
  allowed.
\item \emph{The fleet is unlimited.}  \Cref{def:solution} places no bound on
  $|\Routes|$; by disjointness and non-emptiness, $|\Routes| \le n$ always.
\end{enumerate}

\begin{remark}[Routes versus route sets]
\label{rem:seq}
Two distinct routes may have the same customer set but different visiting
sequences, and \cref{def:route} treats them as different.  This matters: of the
five objectives defined next, only $L(\cdot)$ -- hence $f_4$ and $f_5$ -- is
sensitive to the visiting order.  \Cref{prop:degeneracy} makes the consequences
precise.
\end{remark}

% =========================================================================
\section{The objectives}
\label{sec:objectives}

Fix a solution $\Routes \in \Feas$.  When $\Routes \ne \emptyset$, write
\begin{equation}
  \Lmax = \max_{r \in \Routes} L(r), \quad
  \Wmax = \max_{r \in \Routes} W(r), \quad
  \Omax = \max_{r \in \Routes} O(r),
  \label{eq:maxima}
\end{equation}
and likewise $L^{\min}, W^{\min}, O^{\min}$ for the corresponding minima.

\subsection{Definition}

\begin{definition}[Objective vector]
\label{def:objectives}
The five objectives of \textup{\textsc{MO-CVRP-OS}} are
\begin{align}
  f_1(\Routes) &\;=\; \sum_{i \in \Ground(\Routes)} o_i
    &&\longrightarrow\; \max
    &&\text{(delivered orders)},
  \label{eq:f1}\\[2pt]
  f_2(\Routes) &\;=\; |\Routes|
    &&\longrightarrow\; \min
    &&\text{(routes used)},
  \label{eq:f2}\\[2pt]
  f_3(\Routes) &\;=\; \max_{r \in \Routes} D(r)
    &&\longrightarrow\; \min
    &&\text{(largest route diameter)},
  \label{eq:f3}\\[2pt]
  f_4(\Routes) &\;=\; \min_{r \in \Routes} \beta(r)
    &&\longrightarrow\; \max
    &&\text{(route balance)},
  \label{eq:f4}\\[2pt]
  f_5(\Routes) &\;=\; \sum_{r \in \Routes} L(r)
    &&\longrightarrow\; \min
    &&\text{(total distance)},
  \label{eq:f5}
\end{align}
where the \emph{balance} of a route $r \in \Routes$ is
\begin{equation}
  \beta(r) \;=\; \min\left\{
    \frac{L(r)}{\Lmax}, \;
    \frac{W(r)/Q}{\Wmax/Q}, \;
    \frac{O(r)}{\Omax}
  \right\}
  \;=\;
  \min\left\{
    \frac{L(r)}{\Lmax}, \;
    \frac{W(r)}{\Wmax}, \;
    \frac{O(r)}{\Omax}
  \right\}.
  \label{eq:beta}
\end{equation}
For the empty solution, $f_1(\emptyset) = f_2(\emptyset) = f_3(\emptyset) =
f_5(\emptyset) = 0$ and $f_4(\emptyset) = 1$ by convention.
\end{definition}

The three components of \eqref{eq:beta} are exactly the three normalisations
required: the route length divided by the longest route length, the capacity
utilisation $W(r)/Q$ divided by the largest capacity utilisation among the
routes -- in which the common factor $1/Q$ cancels, so that the component is
simply the load ratio -- and the number of delivered orders divided by the
largest number of orders on any route.  Each component lies in $[0,1]$ and
equals $1$ precisely for a route attaining the corresponding maximum.

The combination rule in \eqref{eq:beta} is the \emph{egalitarian} one: the
balance of a route is its worst normalised component, and the balance of the
solution is the balance of its worst route.  Since a minimum of minima is a
minimum over the union,
\begin{equation}
  f_4(\Routes)
  \;=\;
  \min\left\{
    \frac{L^{\min}}{\Lmax}, \;
    \frac{W^{\min}}{\Wmax}, \;
    \frac{O^{\min}}{\Omax}
  \right\},
  \label{eq:f4-collapsed}
\end{equation}
so $f_4$ is the worst of the three \emph{spread ratios} of the solution.  This
identity is what makes the linearisation of \cref{sec:milp-balance} compact: a
single scalar variable suffices, instead of one ratio variable per route.
Weighted and per-route composite alternatives are discussed in
\cref{sec:alt-balance}.

\subsection{Well-definedness and elementary properties}

\begin{lemma}[The denominators are positive]
\label{lem:welldef}
For every non-empty solution $\Routes \in \Feas$:
$\Lmax \ge L^{\min} \ge 2$, $\Wmax \ge W^{\min} \ge 1$ and
$\Omax \ge O^{\min} \ge 1$.  Consequently \eqref{eq:beta} and
\eqref{eq:f4-collapsed} are well defined, with no division by zero and no need
for a safeguard convention.
\end{lemma}

\begin{proof}
Let $r = (0,i_1,\dots,i_m,0) \in \Routes$ with $m \ge 1$.  All terms of
\eqref{eq:Lr} are non-negative, and the first and last are $d_{0 i_1} \ge 1$
and $d_{i_m 0} \ge 1$ by \cref{as:depot}; if $m = 1$ these are the same two
terms of $L(r) = 2 d_{0 i_1} \ge 2$, and if $m \ge 2$ they are distinct terms,
so $L(r) \ge 2$ in both cases.  Similarly $W(r) \ge q_{i_1} \ge 1$ by
\cref{as:demand} and $O(r) \ge o_{i_1} \ge 1$ by \cref{as:orders}.
\end{proof}

\begin{lemma}[Range and tightness of $f_4$]
\label{lem:f4-range}
For every $\Routes \in \Feas$, $f_4(\Routes) \in (0,1]$.  Moreover
$f_4(\Routes) = 1$ if and only if all routes of $\Routes$ have the same length,
the same load and the same order count; in particular $f_4(\Routes) = 1$
whenever $|\Routes| \le 1$.
\end{lemma}

\begin{proof}
The case $\Routes = \emptyset$ holds by convention.  Otherwise each component
of \eqref{eq:f4-collapsed} is a ratio of two positive numbers
(\cref{lem:welldef}) in which the numerator is a minimum and the denominator
the maximum of the same finite family, hence lies in $(0,1]$; the minimum of
three such numbers lies in $(0,1]$.  Equality with $1$ holds iff each of the
three ratios equals $1$, i.e.\ iff $L^{\min} = \Lmax$, $W^{\min} = \Wmax$ and
$O^{\min} = \Omax$, i.e.\ iff the three quantities are constant over
$\Routes$.  If $|\Routes| = 1$ this is vacuous.
\end{proof}

\begin{lemma}[Integrality]
\label{lem:integrality}
$f_1$, $f_2$, $f_3$ and $f_5$ take values in $\Z_{\ge 0}$, and $f_4$ takes
values in $\mathbb{Q} \cap (0,1]$ with denominator at most
$\max\{\Lmax, \Wmax, \Omax\}$.
\end{lemma}

\begin{proof}
Immediate from \cref{as:dist,as:demand,as:orders}: all of $d$, $q$ and $o$ are
integer valued, so $L(\cdot)$, $W(\cdot)$, $O(\cdot)$ and $D(\cdot)$ are, and
$f_4$ is a ratio of two such integers.
\end{proof}

\subsection{Why a fifth objective}

The four objectives $f_1,\dots,f_4$ are the ones of interest: service level,
fleet size, geographic compactness and workload fairness.  They have, however,
a structural blind spot.

\begin{proposition}[Length degeneracy of $f_1,\dots,f_4$]
\label{prop:degeneracy}
Let $\Routes \in \Feas$ and let $\Routes'$ be obtained from $\Routes$ by
permuting the visiting sequence inside one or more routes, leaving all customer
sets unchanged.  Then
\begin{equation}
  f_1(\Routes') = f_1(\Routes), \qquad
  f_2(\Routes') = f_2(\Routes), \qquad
  f_3(\Routes') = f_3(\Routes),
\end{equation}
while $f_4(\Routes')$ and $f_5(\Routes')$ may differ from $f_4(\Routes)$ and
$f_5(\Routes)$.  Furthermore, the length ratio in \eqref{eq:beta} is invariant
under a common rescaling of all route lengths, so nothing in $f_1,\dots,f_4$
penalises long routes as such: a solution $\Routes'$ obtained from $\Routes$ by
lengthening its shorter routes until all route lengths agree satisfies
$f_4(\Routes') \ge f_4(\Routes)$ and $f_5(\Routes') > f_5(\Routes)$.
Consequently, under the objective vector
$(f_1,f_2,f_3,f_4)$ alone, a solution may be dominated by one obtained from it
purely by \emph{lengthening} routes.
\end{proposition}

\begin{proof}
$f_1$ depends on $\Routes$ only through $\Ground(\Routes)$, $f_2$ only through
$|\Routes|$, and $f_3$ only through the multiset $\{C(r)\}_{r\in\Routes}$ via
\eqref{eq:Dr}; none of the three is affected by the visiting order.  The same
holds for the load and order components of \eqref{eq:beta}, by \eqref{eq:Wr}
and \eqref{eq:Or}.  Only the component $L(r)/\Lmax$ sees the sequence, and it
is a \emph{ratio}: it rewards routes of \emph{equal} length, not of small
length.  \Cref{ex:degeneracy} exhibits a concrete pair $\Routes, \Routes'$ with
$f_4(\Routes') = 1 > f_4(\Routes)$, $f_5(\Routes') > f_5(\Routes)$ and all
other objectives equal.
\end{proof}

\begin{example}[Levelling down]
\label{ex:degeneracy}
Take $n = 6$, $Q = 15$, $q_i = 5$ and $o_i = 1$ for all $i$, with
\texttt{EUC\_2D} coordinates
\[
  \text{depot } 0 = (0,0), \quad
  1 = (1,0),\; 2 = (2,0),\; 3 = (3,0), \quad
  4 = (0,2),\; 5 = (0,3),\; 6 = (0,4),
\]
so that $d$ is integer valued by \eqref{eq:euc2d}.  Consider the two solutions
\[
  \Routes = \bigl\{ (0,1,2,3,0),\; (0,4,5,6,0) \bigr\},
  \qquad
  \Routes' = \bigl\{ (0,2,1,3,0),\; (0,4,5,6,0) \bigr\},
\]
which serve the same customers with the same two customer sets
$\{1,2,3\}$ and $\{4,5,6\}$, each of load $15 = Q$.  Their route lengths are
\begin{gather*}
  L(0,1,2,3,0) = 1+1+1+3 = 6, \qquad
  L(0,2,1,3,0) = 2+1+2+3 = 8, \\
  L(0,4,5,6,0) = 2+1+1+4 = 8 .
\end{gather*}
Both solutions have $f_1 = 6$, $f_2 = 2$ and $f_3 = \max\{d_{13}, d_{46}\} =
\max\{2,2\} = 2$, and both have all loads and order counts equal, so the load
and order components of \eqref{eq:f4-collapsed} equal $1$ in both.  Hence
\[
  f_4(\Routes) = \tfrac{6}{8} = 0.75, \quad f_5(\Routes) = 14,
  \qquad
  f_4(\Routes') = \tfrac{8}{8} = 1, \quad f_5(\Routes') = 16 .
\]
Under $(f_1,f_2,f_3,f_4)$, $\Routes'$ \emph{dominates} $\Routes$: perfect
balance has been achieved by driving two extra kilometres for no reason.  Under
$(f_1,\dots,f_5)$ the two are incomparable, as they should be, and the
trade-off is made explicit to the decision maker.
\end{example}

This phenomenon -- an equity criterion improved by making the best-off worse
rather than the worst-off better -- is well documented in the literature on
workload equity in vehicle routing \cite{Matl2018}, and is the reason why the
total distance $f_5$ is carried as a fifth objective rather than dropped.
\Cref{sec:alt-f5} records the alternatives (a lexicographic tie-break, an
$\varepsilon$-constraint on $f_5$, or restricting routes to be
sequence-optimal), and how to recover the pure four-objective problem.

\begin{remark}[Trivial efficient points]
\label{rem:trivial}
The empty solution attains the optimum of $f_2$, $f_3$, $f_4$ and $f_5$
simultaneously and is therefore always efficient, at the price of the worst
possible $f_1$.  More generally, every single-route solution has
$f_4 = 1$ by \cref{lem:f4-range}.  These points are legitimate but
uninformative; if they are to be excluded, the natural device is a minimum
service level $f_1 \ge \rho \sum_{i\in\Cust} o_i$ imposed as a constraint
(\cref{sec:alt-side}), not a modification of the objectives.
\end{remark}

% =========================================================================
\section{Multi-objective terminology}
\label{sec:pareto}

To keep a single convention, collect the objectives of
\cref{def:objectives} into the \emph{minimisation form}
\begin{equation}
  F(\Routes) \;=\;
  \bigl( -f_1(\Routes),\, f_2(\Routes),\, f_3(\Routes),\, -f_4(\Routes),\,
         f_5(\Routes) \bigr) \;\in\; \R^5 ,
  \label{eq:objvec}
\end{equation}
so that \textup{\textsc{MO-CVRP-OS}} reads $\min_{\Routes \in \Feas} F(\Routes)$ in the
sense of vector optimisation \cite{Ehrgott2005}.

\begin{definition}[Dominance and efficiency]
\label{def:dominance}
For $\Routes, \Routes' \in \Feas$, $\Routes$ \emph{dominates} $\Routes'$,
written $\Routes \prec \Routes'$, if $F_t(\Routes) \le F_t(\Routes')$ for all
$t \in \{1,\dots,5\}$ with at least one strict inequality.  A solution
$\Routes^\star \in \Feas$ is \emph{efficient} (Pareto optimal) if no
$\Routes \in \Feas$ dominates it; the set of efficient solutions is the
\emph{efficient set} $\Feas_E$, and its image $F(\Feas_E)$ is the
\emph{non-dominated set} or \emph{Pareto front}.  An efficient solution is
\emph{supported} if it minimises $\sum_t w_t F_t(\cdot)$ over $\Feas$ for some
weight vector $w \in \R^5_{>0}$, and \emph{unsupported} otherwise.
\end{definition}

\begin{definition}[Ideal and nadir points]
The \emph{ideal point} is $F^{\mathrm{id}}_t = \min_{\Routes \in \Feas}
F_t(\Routes)$ and the \emph{nadir point} is $F^{\mathrm{nad}}_t =
\max_{\Routes \in \Feas_E} F_t(\Routes)$, for $t = 1,\dots,5$.
\end{definition}

Two consequences of \cref{sec:objectives} are worth stating.  First,
$\Feas$ is finite, so $\Feas_E \ne \emptyset$ and the Pareto front is finite;
by \cref{lem:integrality} it is a finite subset of
$\Z^3 \times \mathbb{Q} \times \Z$ up to the sign convention.  Second, $\Feas$
is a set of combinatorial objects and not a convex set, so the Pareto front is
in general non-convex and contains unsupported points; weighted-sum
scalarisation cannot enumerate it, whereas the $\varepsilon$-constraint method
can \cite{Ehrgott2005}.  \Cref{sec:alt-eps} shows that the
$\varepsilon$-constraint treatment of $f_4$ is particularly convenient here,
because bounding $f_4$ from below by a \emph{constant} removes the only
non-linearity of the model.

Part of the ideal point is available in closed form.  Since the empty solution
is feasible, $f_2, f_3, f_5$ attain the value $0$ and $f_4$ attains the value
$1$, so
\begin{equation}
  F^{\mathrm{id}} \;=\;
  \Bigl( -\textstyle\sum_{i \in \Cust} o_i,\; 0,\; 0,\; -1,\; 0 \Bigr),
\end{equation}
where the first component uses that serving every customer is always possible
(put each customer on its own route; feasible by \cref{as:demand}).  The ideal
point is thus never attained by a single solution as soon as $n \ge 1$, which
is the usual situation in multi-objective optimisation and confirms that the
five criteria genuinely conflict.

% =========================================================================
\section{Mixed-integer linear formulation}
\label{sec:milp}

\subsection{Overview and design decisions}

Three of the five objectives are \emph{route-level}: $f_2$ counts routes, $f_3$
is a maximum over routes, and $f_4$ compares routes to one another.  A model
must therefore be able to name an individual route, which rules out the classic
two-index CVRP formulation in which routes are recovered only implicitly from
the arc-incidence vector.  The formulation below is consequently
\emph{vehicle indexed} (three index): a set $\Veh$ of interchangeable vehicles
is introduced, and route $k$ is the tour performed by vehicle $k$.
\Cref{sec:alt-formulations} discusses the two alternatives -- a two-index model
with auxiliary route-identification machinery, and a set-partitioning model in
which routes are explicit columns.

Vehicle indexing costs two things, and both are addressed below: symmetry
among the interchangeable vehicles (\cref{sec:milp-symmetry}) and a factor
$|\Veh|$ in the model size (\cref{sec:milp-size}).

\subsection{Sets, parameters and constants}
\label{sec:milp-data}

\begin{center}
\begin{tabular}{@{}ll@{}}
\toprule
Symbol & Meaning \\
\midrule
$\Cust = \{1,\dots,n\}$ & customers \\
$\Nodes = \{0\} \cup \Cust$ & nodes; $0$ is the depot \\
$\Arcs = \{(i,j) \in \Nodes^2 : i \ne j\}$ & arcs \\
$\Veh = \{1,\dots,\overline{K}\}$ & vehicles (interchangeable, identical) \\
$d_{ij} \in \Z_{\ge0}$ & distance, symmetric, $d_{ii}=0$ (\cref{as:dist}) \\
$q_i \in \Z_{\ge1}$, $i \in \Cust$ & demand, $q_i \le Q$ (\cref{as:demand}) \\
$o_i \in \Z_{\ge1}$, $i \in \Cust$ & orders, default $1$ (\cref{as:orders}) \\
$Q \in \Z_{>0}$ & vehicle capacity \\
$\overline{K}$ & fleet-size bound; default $\overline{K} = n$ \\
$\Lub$ & upper bound on a route length \\
$\Oub = \sum_{i\in\Cust} o_i$ & upper bound on a route order count \\
$P \in \Z_{\ge 1}$ & resolution parameter of the balance linearisation \\
$\Gamma = 2^{P}$ & number of grid steps, see \cref{sec:milp-balance} \\
\bottomrule
\end{tabular}
\end{center}

\paragraph{Fleet-size bound.}
By \cref{def:solution} every route is non-empty and customer sets are disjoint,
so $|\Routes| \le n$ and $\overline{K} = n$ excludes no solution.  Any smaller
$\overline{K}$ is a genuine restriction of $\Feas$, and one that is not
innocuous here: the solutions that minimise $f_3$ are typically those that
split customers into many small routes -- with $\overline{K} = n$, putting every
customer on its own route gives $f_3 = 0$ -- so truncating $\overline{K}$ cuts
away exactly the part of the Pareto front where $f_3$ is best.  When a smaller
$\overline{K}$ is used for tractability this must be reported as part of the
problem instance.

\paragraph{Length bound.}
$\Lub$ must exceed the length of any single route.  A route uses at most $n+1$
arcs, so
\begin{equation}
  \Lub \;=\; (n+1)\, d^{\max}, \qquad
  d^{\max} = \max_{(i,j) \in \Arcs} d_{ij}
  \label{eq:Lub}
\end{equation}
is valid.  A tighter valid bound, useful because $\Lub$ appears in big-$M$
coefficients, is obtained from the capacity: a route carries load at most $Q$,
so it visits at most $m^{\max} = \max\{ m : \text{the } m \text{ smallest
demands sum to at most } Q \}$ customers, and
$\Lub = (m^{\max}+1)\,d^{\max}$ is valid as well.  Both bounds are computed
from the instance data and require no assumption on $d$.

\subsection{Decision variables}
\label{sec:milp-vars}

\begin{align}
  x^k_{ij} &\in \{0,1\}, & &(i,j) \in \Arcs,\ k \in \Veh
  && \text{vehicle } k \text{ traverses } (i,j), \\
  y^k_i    &\in \{0,1\}, & &i \in \Cust,\ k \in \Veh
  && \text{vehicle } k \text{ serves customer } i, \\
  z_i      &\in \{0,1\}, & &i \in \Cust
  && \text{customer } i \text{ is served at all}, \\
  u^k      &\in \{0,1\}, & &k \in \Veh
  && \text{vehicle } k \text{ performs a route}, \\
  g_{ij}   &\ge 0,       & &(i,j) \in \Arcs
  && \text{load carried on arc } (i,j), \\
  D^k      &\ge 0,       & &k \in \Veh
  && \text{diameter of route } k, \\
  L^k,\, W^k,\, O^k &\ge 0, & &k \in \Veh
  && \text{length, load, order count of route } k, \\
  \Dmax,\, \Lmax,\, \Wmax,\, \Omax &\ge 0, & &
  && \text{the corresponding maxima}, \\
  \lambda  &\in [0,1],   & &
  && \text{route balance}, \\
  b_p      &\in \{0,1\}, & &p \in \{0,\dots,P\}
  && \text{binary expansion of } \lambda, \\
  t^{L}_p,\, t^{W}_p,\, t^{O}_p &\ge 0, & &p \in \{0,\dots,P\}
  && \text{linearised products.}
\end{align}
It is convenient to abbreviate the \emph{aggregated} arc variables
\begin{equation}
  x_{ij} \;=\; \sum_{k \in \Veh} x^k_{ij}, \qquad (i,j) \in \Arcs,
  \label{eq:aggregate}
\end{equation}
which may be substituted out or kept as defined variables.

\subsection{Routing, assignment and capacity}
\label{sec:milp-routing}

\begin{align}
  \sum_{k \in \Veh} y^k_i &\;=\; z_i,
  & &i \in \Cust, \label{eq:c-assign}\\
  \sum_{j \in \Nodes \setminus \{i\}} x^k_{ij} &\;=\; y^k_i,
  & &i \in \Cust,\ k \in \Veh, \label{eq:c-out}\\
  \sum_{j \in \Nodes \setminus \{i\}} x^k_{ji} &\;=\; y^k_i,
  & &i \in \Cust,\ k \in \Veh, \label{eq:c-in}\\
  \sum_{j \in \Cust} x^k_{0j} \;=\; \sum_{j \in \Cust} x^k_{j0}
  &\;=\; u^k,
  & &k \in \Veh, \label{eq:c-depot}\\
  \sum_{i \in \Cust} q_i\, y^k_i &\;\le\; Q\, u^k,
  & &k \in \Veh, \label{eq:c-cap}\\
  u^k &\;\le\; \sum_{i \in \Cust} y^k_i,
  & &k \in \Veh. \label{eq:c-nonempty}
\end{align}

Constraint~\eqref{eq:c-assign} is the optional-service constraint: a customer is
assigned to at most one vehicle, and $z_i$ records whether it is assigned at
all.  Replacing ``$= z_i$'' by ``$= 1$'' recovers the mandatory-service CVRP.
Constraints \eqref{eq:c-out}--\eqref{eq:c-in} are the degree constraints,
written so that a customer has in-degree and out-degree $1$ in the
sub-graph of vehicle $k$ exactly when $k$ serves it.  \Cref{eq:c-depot} lets
each used vehicle leave and re-enter the depot exactly once, which forbids
multiple trips per vehicle.  \Cref{eq:c-cap} is the capacity constraint and
simultaneously forces $y^k \equiv 0$ when $u^k = 0$, because $q_i \ge 1$ for
all $i$ (\cref{as:demand}); the explicit linking inequalities
$y^k_i \le u^k$ are therefore redundant, though they strengthen the linear
relaxation and may be added.  \Cref{eq:c-nonempty} forbids a ``used'' vehicle
with no customers, so that $\sum_k u^k$ counts routes in the sense of
\cref{def:route}.

\subsection{Connectivity}
\label{sec:milp-flow}

Constraints \eqref{eq:c-out}--\eqref{eq:c-depot} force the arc set of each
vehicle to be a disjoint union of circuits, but do not yet force those circuits
to pass through the depot.  Connectivity is imposed by a single-commodity flow
of Gavish--Graves type \cite{GavishGraves1978} on the aggregated arcs
\eqref{eq:aggregate}, where $g_{ij}$ is interpreted as the load still on board
when arc $(i,j)$ is traversed:
\begin{align}
  \sum_{j \in \Nodes \setminus \{i\}} g_{ji}
  \;-\; \sum_{j \in \Nodes \setminus \{i\}} g_{ij}
  &\;=\; q_i\, z_i,
  & &i \in \Cust, \label{eq:c-flowcons}\\
  q_j\, x_{0j} \;\le\; g_{0j} &\;\le\; Q\, x_{0j},
  & &j \in \Cust, \label{eq:c-flowdepot}\\
  q_j\, x_{ij} \;\le\; g_{ij} &\;\le\; (Q - q_i)\, x_{ij},
  & &i,j \in \Cust,\ i \ne j, \label{eq:c-flowcust}\\
  g_{i0} &\;=\; 0,
  & &i \in \Cust. \label{eq:c-flowreturn}
\end{align}
\Cref{eq:c-flowcons} states that each served customer absorbs its demand from
the flow, and each unserved customer absorbs nothing.  \Cref{eq:c-flowdepot}
bounds the load leaving the depot by the capacity, \cref{eq:c-flowcust} states
that after visiting $i$ the vehicle carries at most $Q - q_i$ and, if it goes
on to $j$, at least $q_j$, and \cref{eq:c-flowreturn} empties the vehicle on
its return leg.  The lower bounds in
\eqref{eq:c-flowdepot}--\eqref{eq:c-flowcust} are valid and tighten the linear
relaxation; they are not needed for correctness.

\subsection{Symmetry breaking}
\label{sec:milp-symmetry}

The vehicles are interchangeable, so every solution has up to
$\overline{K}!/(\overline{K}-|\Routes|)!$ encodings.  The following two
families of constraints select the canonical one in which routes are ordered by
increasing lowest-index customer:
\begin{align}
  u^{k} &\;\ge\; u^{k+1},
  & &k \in \{1,\dots,\overline{K}-1\}, \label{eq:c-sym1}\\
  y^k_i &\;\le\; \sum_{j \in \Cust,\; j < i} y^{k-1}_j,
  & &i \in \Cust,\ k \in \Veh \setminus \{1\}. \label{eq:c-sym2}
\end{align}

\begin{lemma}
\label{lem:symmetry}
\Cref{eq:c-sym1,eq:c-sym2} cut off no solution of $\Feas$ with at most
$\overline{K}$ routes: every such solution admits an encoding satisfying them.
\end{lemma}

\begin{proof}
Given $\Routes$ with $|\Routes| = \kappa \le \overline{K}$, order its routes as
$r_1,\dots,r_\kappa$ by increasing $\mu(r) = \min C(r)$; the $\mu$-values are
distinct because the customer sets are disjoint and non-empty.  Assign $r_k$ to
vehicle $k$ for $k \le \kappa$ and set $u^k = 1$ iff $k \le \kappa$, which
satisfies \eqref{eq:c-sym1}.  For \eqref{eq:c-sym2}, let $k \ge 2$ and
$i \in \Cust$ with $y^k_i = 1$.  Then $k \le \kappa$ and
$i \in C(r_k)$, so $i \ge \mu(r_k) > \mu(r_{k-1})$, and the customer
$j = \mu(r_{k-1}) < i$ satisfies $y^{k-1}_j = 1$, so the right-hand side is at
least $1$.  If $y^k_i = 0$ the constraint is trivially satisfied because its
right-hand side is non-negative.
\end{proof}

\subsection{Route attributes and the diameter}
\label{sec:milp-attributes}

\begin{align}
  L^k &\;=\; \sum_{(i,j) \in \Arcs} d_{ij}\, x^k_{ij},
  & &k \in \Veh, \label{eq:c-Lk}\\
  W^k &\;=\; \sum_{i \in \Cust} q_i\, y^k_i,
  & &k \in \Veh, \label{eq:c-Wk}\\
  O^k &\;=\; \sum_{i \in \Cust} o_i\, y^k_i,
  & &k \in \Veh, \label{eq:c-Ok}\\
  D^k &\;\ge\; d_{ij}\,\bigl( y^k_i + y^k_j - 1 \bigr),
  & &i,j \in \Cust,\ i < j,\ k \in \Veh, \label{eq:c-Dk}\\
  \Dmax &\;\ge\; D^k,
  & &k \in \Veh. \label{eq:c-Dmax}
\end{align}

\Cref{eq:c-Dk} is the standard product linearisation for the term
$d_{ij}\, y^k_i y^k_j$: it is active only when both customers are on route $k$,
in which case it reads $D^k \ge d_{ij}$, and is implied by $D^k \ge 0$
otherwise since $d_{ij} \ge 0$.  Only unordered pairs are needed because $d$ is
symmetric (\cref{as:dist}).  Note that \eqref{eq:c-Dk} involves the assignment
variables only, which is \cref{prop:degeneracy} in constraint form.

\subsection{Route balance}
\label{sec:milp-balance}

The maxima are bounded from below, one-sidedly:
\begin{align}
  \Lmax \;\ge\; L^k, \qquad
  \Wmax \;\ge\; W^k, \qquad
  \Omax \;\ge\; O^k,
  \qquad & k \in \Veh, \label{eq:c-maxima}\\
  \Lmax \;\le\; \Lub, \qquad
  \Wmax \;\le\; Q, \qquad
  \Omax \;\le\; \Oub. &  \label{eq:c-maxima-ub}
\end{align}
The upper bounds \eqref{eq:c-maxima-ub} are valid because every individual
route satisfies $L^k \le \Lub$, $W^k \le Q$ and $O^k \le \Oub$; they are stated
explicitly because the product linearisation of \cref{lem:mccormick} requires a
finite upper bound on each factor.
The balance variable $\lambda$ is then bounded above by each of the three
normalised components of \eqref{eq:beta}, for used routes only.  Written
directly, this is
\begin{align}
  \lambda \cdot \Lmax &\;\le\; L^k + \Lub\,(1-u^k), & &k \in \Veh,
  \label{eq:c-balL}\\
  \lambda \cdot \Wmax &\;\le\; W^k + Q\,(1-u^k), & &k \in \Veh,
  \label{eq:c-balW}\\
  \lambda \cdot \Omax &\;\le\; O^k + \Oub\,(1-u^k), & &k \in \Veh.
  \label{eq:c-balO}
\end{align}
The big-$M$ coefficients are the smallest valid ones: the left-hand sides are
bounded by $\Lmax \le \Lub$, $\Wmax \le Q$ and $\Omax \le \Oub$ respectively,
since $\lambda \le 1$.  An unused vehicle ($u^k = 0$) is thereby excluded from
the minimum, as \cref{eq:f4} requires.

Constraints \eqref{eq:c-balL}--\eqref{eq:c-balO} are the only non-linear part
of the model: each contains the product of the continuous variable $\lambda$
with a continuous variable.  They are linearised exactly by restricting
$\lambda$ to a uniform grid of $\Gamma + 1 = 2^P + 1$ points in $[0,1]$ and
expanding the grid index in binary:
\begin{align}
  \lambda &\;=\; \frac{1}{\Gamma} \sum_{p=0}^{P} 2^{p}\, b_p,
  \qquad
  \sum_{p=0}^{P} 2^{p}\, b_p \;\le\; \Gamma,
  \label{eq:c-lambda-exp}
\end{align}
together with, for each $p \in \{0,\dots,P\}$, the McCormick identities for the
products $t^L_p = b_p \Lmax$, $t^W_p = b_p \Wmax$ and $t^O_p = b_p \Omax$:
\begin{align}
  t^{L}_p \;\le\; \Lub\, b_p, \quad
  t^{L}_p \;\le\; \Lmax, \quad
  t^{L}_p \;\ge\; \Lmax - \Lub\,(1 - b_p), \quad
  t^{L}_p \;\ge\; 0,
  \label{eq:c-mcL}
\end{align}
and the analogous groups for $t^W_p$ (with $\Wmax$, $Q$) and $t^O_p$ (with
$\Omax$, $\Oub$).  Substituting
$\lambda \Lmax = \Gamma^{-1} \sum_p 2^p t^L_p$ and its analogues,
\cref{eq:c-balL,eq:c-balW,eq:c-balO} become the linear constraints
\begin{align}
  \frac{1}{\Gamma} \sum_{p=0}^{P} 2^{p}\, t^{L}_p
    &\;\le\; L^k + \Lub\,(1-u^k), & &k \in \Veh,
  \label{eq:c-balL-lin}\\
  \frac{1}{\Gamma} \sum_{p=0}^{P} 2^{p}\, t^{W}_p
    &\;\le\; W^k + Q\,(1-u^k), & &k \in \Veh,
  \label{eq:c-balW-lin}\\
  \frac{1}{\Gamma} \sum_{p=0}^{P} 2^{p}\, t^{O}_p
    &\;\le\; O^k + \Oub\,(1-u^k), & &k \in \Veh.
  \label{eq:c-balO-lin}
\end{align}

\subsection{The model}
\label{sec:milp-model}

\begin{definition}[\textup{\textsc{MO-CVRP-OS}} MILP]
\label{def:milp}
\begin{align}
  \text{maximise } \quad & f_1 = \sum_{i \in \Cust} o_i\, z_i, \label{eq:o1}\\
  \text{minimise } \quad & f_2 = \sum_{k \in \Veh} u^k, \label{eq:o2}\\
  \text{minimise } \quad & f_3 = \Dmax, \label{eq:o3}\\
  \text{maximise } \quad & f_4 = \lambda, \label{eq:o4}\\
  \text{minimise } \quad & f_5 = \sum_{k \in \Veh} L^k
                               = \sum_{(i,j) \in \Arcs} d_{ij}\, x_{ij},
  \label{eq:o5}
\end{align}
subject to
\begin{center}
\begin{tabular}{@{}ll@{}}
  \eqref{eq:c-assign}--\eqref{eq:c-nonempty}
    & assignment, degree, depot and capacity constraints, \\
  \eqref{eq:c-flowcons}--\eqref{eq:c-flowreturn}
    & connectivity, \\
  \eqref{eq:c-sym1}--\eqref{eq:c-sym2}
    & symmetry breaking, \\
  \eqref{eq:c-Lk}--\eqref{eq:c-Dmax}
    & route attributes and diameter, \\
  \eqref{eq:c-maxima}--\eqref{eq:c-balO-lin}
    & route balance, \\
\end{tabular}
\end{center}
and to the variable restrictions
\begin{equation}
  x,\,y,\,z,\,u,\,b \ \text{binary}, \qquad
  g,\,D,\,L,\,W,\,O,\,\Dmax,\,\Lmax,\,\Wmax,\,\Omax,\,t \;\ge\; 0, \qquad
  \lambda \in [0,1].
  \label{eq:o-domains}
\end{equation}
\end{definition}

\subsection{Correctness}
\label{sec:milp-correct}

\begin{proposition}[The feasible sets correspond]
\label{prop:validity}
Let $\overline{K} \ge 1$.
\begin{enumerate}[label=(\alph*),leftmargin=*]
\item Every $(x,y,z,u,g)$ satisfying
  \eqref{eq:c-assign}--\eqref{eq:c-flowreturn} encodes a solution
  $\Routes \in \Feas$ with $|\Routes| = \sum_{k} u^k \le \overline{K}$: for each
  $k$ with $u^k = 1$ the arcs $\{(i,j) : x^k_{ij} = 1\}$ form a single circuit
  through the depot, whose customer set is $\{i : y^k_i = 1\}$ and whose load is
  at most $Q$; for each $k$ with $u^k = 0$ that arc set is empty.
\item Conversely, every $\Routes \in \Feas$ with $|\Routes| \le \overline{K}$
  admits an encoding satisfying \eqref{eq:c-assign}--\eqref{eq:c-flowreturn}
  and the symmetry-breaking constraints \eqref{eq:c-sym1}--\eqref{eq:c-sym2}.
\end{enumerate}
\end{proposition}

\begin{proof}
(a) Fix $k$.  If $u^k = 0$ then \eqref{eq:c-cap} gives
$\sum_i q_i y^k_i \le 0$, hence $y^k \equiv 0$ by \cref{as:demand}, hence
$x^k \equiv 0$ by \eqref{eq:c-out} and \eqref{eq:c-depot}.  If $u^k = 1$, then
by \eqref{eq:c-out}--\eqref{eq:c-depot} every node of
$\{0\} \cup C_k$, where $C_k = \{i : y^k_i = 1\}$, has out-degree and in-degree
exactly $1$ in the arc set of vehicle $k$, and every other node has degree $0$;
such an arc set is a disjoint union of circuits covering $\{0\} \cup C_k$, and
$C_k \ne \emptyset$ by \eqref{eq:c-nonempty}.

It remains to exclude circuits avoiding the depot.  Summing \eqref{eq:c-assign}
over $k$ shows that in the aggregated digraph
$\widehat{G} = \bigl(\Nodes, \{(i,j) : x_{ij} \ge 1\}\bigr)$ every customer $i$ has out-degree and
in-degree exactly $z_i \in \{0,1\}$; in particular $x_{ij} \in \{0,1\}$ for all
arcs, and the arcs of the different vehicles are pairwise disjoint.  Suppose
some vehicle's arc set contained a circuit $Z$ with $V(Z) \subseteq \Cust$.
Every $i \in V(Z)$ has $z_i = 1$ and its unique incoming and unique outgoing
arc in $\widehat{G}$ both belong to $Z$; therefore every arc of
$\Arcs$ with exactly one endpoint in $V(Z)$ has $x_{ij} = 0$, and by
\eqref{eq:c-flowdepot}--\eqref{eq:c-flowcust} carries $g_{ij} = 0$.  Summing
\eqref{eq:c-flowcons} over $i \in V(Z)$, the flow on the arcs internal to $Z$
appears once positively and once negatively and cancels, and the arcs leaving
$V(Z)$ carry no flow, so the left-hand side vanishes, while the right-hand side
equals $\sum_{i \in V(Z)} q_i \ge |V(Z)| \ge 1 > 0$ by \cref{as:demand}.  This
contradiction shows that each vehicle's arc set is a single circuit through the
depot, i.e.\ a route in the sense of \cref{def:route}, with customer set $C_k$.
The customer sets are disjoint by \eqref{eq:c-assign} and the load bound is
\eqref{eq:c-cap}.

(b) Order the routes of $\Routes$ as in the proof of \cref{lem:symmetry} and
assign route $r_k$ to vehicle $k$.  Set $x^k_{ij} = 1$ for the arcs of $r_k$,
$y^k_i = 1$ for $i \in C(r_k)$, $z_i = 1$ for $i \in \Ground(\Routes)$,
$u^k = 1$ for $k \le |\Routes|$, and all remaining variables of these families
to $0$.  Constraints
\eqref{eq:c-assign}--\eqref{eq:c-nonempty} hold by construction and capacity
feasibility.  For the flow, let $g_{ij}$ be the load still on board on arc
$(i,j)$ of the route containing it, i.e.\ for
$r_k = (0,i_1,\dots,i_m,0)$ set
$g_{i_h i_{h+1}} = \sum_{l > h} q_{i_l}$ for $h = 0,\dots,m$ with the
conventions $i_0 = i_{m+1} = 0$, and $g_{ij} = 0$ on all other arcs.  Then
\eqref{eq:c-flowcons} holds at every customer of a route because consecutive
values differ by exactly $q_{i_h}$, and at every unserved customer because all
incident $g$ vanish; \eqref{eq:c-flowdepot} holds because
$g_{0 i_1} = W(r_k) \le Q$ and $g_{0i_1} \ge q_{i_1}$;
\eqref{eq:c-flowcust} holds because
$g_{i_h i_{h+1}} = W(r_k) - \sum_{l \le h} q_{i_l} \le Q - q_{i_h}$ and
$g_{i_h i_{h+1}} \ge q_{i_{h+1}}$ for $h < m$; and \eqref{eq:c-flowreturn}
holds because $g_{i_m 0} = 0$.  The symmetry-breaking constraints hold by
\cref{lem:symmetry}.
\end{proof}

\begin{remark}[Where \cref{as:demand} is used]
\label{rem:demand-needed}
The proof of \cref{prop:validity}(a) uses $q_i \ge 1$ twice: to deduce
$y^k \equiv 0$ from $u^k = 0$, and -- essentially -- to derive the
contradiction $0 = \sum_{i \in V(Z)} q_i > 0$.  If an instance admitted a
customer with zero demand, a depot-free circuit consisting entirely of such
customers would satisfy \eqref{eq:c-flowcons}--\eqref{eq:c-flowreturn} and the
formulation would be wrong.  No instance in this repository has a zero-demand
customer, but a formulation intended for arbitrary data should either add the
usual $\varepsilon$-demand perturbation, or replace the flow system by
generalised subtour elimination constraints (\cref{sec:alt-sec}).
\end{remark}

\begin{proposition}[The auxiliary variables report the intended values]
\label{prop:tightness}
Let $(x,y,z,u,g,\dots)$ be feasible for \cref{def:milp} and let $\Routes$ be the
solution it encodes (\cref{prop:validity}).  Then
\begin{equation}
  \sum_{i \in \Cust} o_i z_i = f_1(\Routes), \quad
  \sum_{k \in \Veh} u^k = f_2(\Routes), \quad
  \sum_{k} L^k = f_5(\Routes),
  \label{eq:exact-objs}
\end{equation}
identically, while
\begin{equation}
  \Dmax \;\ge\; f_3(\Routes)
  \qquad\text{and}\qquad
  \lambda \;\le\; \bigl\lfloor f_4(\Routes) \bigr\rfloor_{\Gamma},
  \label{eq:bound-objs}
\end{equation}
where $\lfloor \cdot \rfloor_{\Gamma}$ rounds down to the grid
$\{0, 1/\Gamma, \dots, 1\}$.  Both bounds in \eqref{eq:bound-objs} are attained
by some feasible completion of the same $(x,y,z,u,g)$.  Consequently, for any
scalarisation that is non-decreasing in $\Dmax$ and non-increasing in
$\lambda$ -- weighted sums with non-negative weights, $\varepsilon$-constraints,
lexicographic orderings, and any combination of these -- an optimal solution of
\cref{def:milp} reports $\Dmax = f_3(\Routes)$ and
$\lambda = \lfloor f_4(\Routes) \rfloor_{\Gamma}$.
\end{proposition}

\begin{proof}
The identities \eqref{eq:exact-objs} are immediate from
\eqref{eq:c-assign}, \eqref{eq:c-nonempty} and \eqref{eq:c-Lk}, the last
because the arc sets of distinct vehicles are disjoint.

For $\Dmax$: by \eqref{eq:c-Dk}, $D^k \ge d_{ij}$ for every pair
$i,j \in C_k$, so $D^k \ge D(r_k)$ and $\Dmax \ge \max_k D(r_k) = f_3(\Routes)$
by \eqref{eq:c-Dmax}; setting $D^k = D(r_k)$ and $\Dmax = f_3(\Routes)$ is
feasible, since \eqref{eq:c-Dk} is then satisfied with equality for a maximising
pair and slack elsewhere, and \eqref{eq:c-Dmax} holds by definition of a
maximum.  As $\Dmax$ occurs in no other constraint and is minimised, it equals
$f_3(\Routes)$ at any optimum.

For $\lambda$: by \eqref{eq:c-maxima}, $\Lmax \ge \max_k L^k = \max_{r \in \Routes} L(r)$ in the
notation of \eqref{eq:maxima}, and similarly for $\Wmax, \Omax$.  For any used
route $k$, \eqref{eq:c-balL} gives
$\lambda \le L^k/\Lmax \le L^k/\max_{r\in\Routes} L(r)$ and analogously for the other two
families, so $\lambda \le \beta(r_k)$ for every $k$ with $u^k = 1$, whence
$\lambda \le f_4(\Routes)$; since \eqref{eq:c-lambda-exp} confines $\lambda$ to
the grid, $\lambda \le \lfloor f_4(\Routes) \rfloor_\Gamma$.  Conversely,
setting $\Lmax = \max_{r} L(r)$, $\Wmax = \max_{r} W(r)$, $\Omax = \max_{r} O(r)$,
$\lambda = \lfloor f_4(\Routes) \rfloor_\Gamma$ with the corresponding
$b$, and $t^L_p = b_p \Lmax$, $t^W_p = b_p \Wmax$, $t^O_p = b_p \Omax$ is
feasible: \eqref{eq:c-mcL} and its analogues hold with equality by
\cref{lem:mccormick} below, and
\eqref{eq:c-balL-lin}--\eqref{eq:c-balO-lin} reduce to
$\lambda \max_r L(r) \le L^k$, $\lambda \max_r W(r) \le W^k$ and
$\lambda \max_r O(r) \le O^k$ for used $k$, which hold because
$\lambda \le f_4(\Routes) \le \beta(r_k)$.  Unused $k$ are covered by the
big-$M$ terms, which are valid by
$\lambda \Lmax \le \Lub$, $\lambda \Wmax \le Q$, $\lambda \Omax \le \Oub$.
Since $\Lmax, \Wmax, \Omax$ occur only in \eqref{eq:c-maxima} and in
\eqref{eq:c-balL}--\eqref{eq:c-balO}, where decreasing them relaxes the
constraint, and $\lambda$ is maximised, the stated values are attained.
\end{proof}

\begin{lemma}[Exactness of the product linearisation]
\label{lem:mccormick}
Let $b \in \{0,1\}$, let $v$ be a variable with $0 \le v \le \overline{v}$, and
let $t$ satisfy $t \le \overline{v} b$, $t \le v$,
$t \ge v - \overline{v}(1-b)$ and $t \ge 0$.  Then $t = bv$.  Consequently, with
\eqref{eq:c-lambda-exp} and \eqref{eq:c-mcL},
$\Gamma^{-1}\sum_{p} 2^{p} t^{L}_p = \lambda \Lmax$ exactly, and analogously
for $W$ and $O$; the constraints
\eqref{eq:c-balL-lin}--\eqref{eq:c-balO-lin} are therefore equivalent to
\eqref{eq:c-balL}--\eqref{eq:c-balO} restricted to $\lambda$ on the grid.
\end{lemma}

\begin{proof}
If $b = 0$ the first and fourth inequalities give $t = 0 = bv$.  If $b = 1$ the
second and third give $v \le t \le v$, so $t = v = bv$.  Summing,
$\Gamma^{-1}\sum_p 2^p t^L_p = \Gamma^{-1}\sum_p 2^p b_p \Lmax
= \lambda \Lmax$ by \eqref{eq:c-lambda-exp}.
\end{proof}

\begin{corollary}[Resolution and exact recovery]
\label{cor:resolution}
Let $M = \max\{\Lub,\, Q,\, \Oub\}$, let $\Routes$ be the solution encoded by an
optimal solution of \cref{def:milp} and let $\lambda$ be the reported balance.
\begin{enumerate}[label=(\alph*),leftmargin=*]
\item $0 \le f_4(\Routes) - \lambda < 2^{-P}$.  In particular the model is
  exact for every value of $f_4$ that is a dyadic rational with denominator
  dividing $\Gamma$, and approximate otherwise.
\item Let $\mathbb{Q}_M = \{ a/c \,:\, c \in \{1,\dots,M\},\ a \in
  \{0,\dots,c\} \}$, so that $f_4(\Routes) \in \mathbb{Q}_M$ by
  \cref{lem:integrality}.  If $P \ge 2\log_2 M$, then
  $[\lambda,\, \lambda + 2^{-P})$ contains exactly one element of
  $\mathbb{Q}_M$, namely $f_4(\Routes)$; the exact value is therefore recovered
  from $\lambda$ in post-processing, for instance by a Stern--Brocot search.
\end{enumerate}
\end{corollary}

\begin{proof}
(a) The grid $\{0,1/\Gamma,\dots,1\}$ has mesh $1/\Gamma = 2^{-P}$, and
\cref{prop:tightness} gives $\lambda = \lfloor f_4(\Routes) \rfloor_\Gamma$.
(b) Let $a/c \ne a'/c'$ be elements of $\mathbb{Q}_M$.  Then
$|a/c - a'/c'| = |ac' - a'c|/(cc') \ge 1/(cc') \ge 1/M^2 \ge 2^{-P}$, the last
step by $P \ge 2\log_2 M$.  Two distinct elements of $\mathbb{Q}_M$ inside
$[\lambda, \lambda + 2^{-P})$ would differ by strictly less than $2^{-P}$, a
contradiction, and $f_4(\Routes)$ lies in that interval by~(a).
\end{proof}

\begin{remark}[Exactness without discretisation]
\label{rem:exactness}
\Cref{cor:resolution} is a statement about a genuine approximation: the model of
\cref{def:milp} is an exact MILP for $f_4$ restricted to a grid, not for $f_4$
itself.  Three exact alternatives exist and are compared in
\cref{sec:alt-eps}: fixing $\lambda = \theta$ as a constant, which makes
\eqref{eq:c-balL}--\eqref{eq:c-balO} linear and is what an
$\varepsilon$-constraint scheme does anyway; bisection on $\theta$; and keeping
\eqref{eq:c-balL}--\eqref{eq:c-balO} as bilinear constraints and handing the
model to a global MIQCP solver.  A fourth option keeps the structure of
\eqref{eq:c-lambda-exp}--\eqref{eq:c-balO-lin} but replaces the dyadic grid by
an arbitrary finite set $\Theta \subseteq [0,1]$, selected by binaries
$s_\theta$ with $\sum_{\theta \in \Theta} s_\theta = 1$ and
$\lambda = \sum_{\theta \in \Theta} \theta\, s_\theta$; taking
$\Theta = \mathbb{Q}_M$ as in \cref{cor:resolution} makes the model exact, at
the price of $\Theta(M^2)$ binaries instead of $P+1$.
\end{remark}

\subsection{Bounds and valid inequalities}
\label{sec:milp-bounds}

\begin{proposition}
\label{prop:bounds}
The following are valid for \cref{def:milp}.
\begin{enumerate}[label=(\alph*),leftmargin=*]
\item \emph{Fleet lower bound.}
  $\sum_{k \in \Veh} u^k \;\ge\; \frac{1}{Q}\sum_{i \in \Cust} q_i z_i$,
  and, when $z$ is fixed, the rounded form
  $\sum_k u^k \ge \bigl\lceil \frac{1}{Q}\sum_i q_i z_i \bigr\rceil$.
\item \emph{Diameter--fleet interaction.}  If $\sum_{k} u^k \le \kappa$ and
  $\sum_{i} z_i = \sigma \ge 1$, then some used route serves at least
  $\lceil \sigma/\kappa \rceil$ customers, so
  \[
    \Dmax \;\ge\;
    \min\Bigl\{ \max_{i,j \in S} d_{ij} \;:\;
      S \subseteq \Cust,\ |S| = \lceil \sigma/\kappa \rceil \Bigr\} .
  \]
  In particular $f_3 = 0$ forces every used route to have pairwise
  zero-distance customers, which by \cref{as:dist} is possible only in the ten
  \textsf{AGS} instances and the twelve \texttt{EXPLICIT} ones, and even there
  only for very small routes.  Without a bound on $\sum_k u^k$ no non-trivial
  lower bound on $\Dmax$ exists, since putting every customer on its own route
  yields $f_3 = 0$.  The bound above is stated for insight: evaluating it is
  itself NP-hard in general.
\item \emph{Balance upper bound.}  $\lambda \le 1$, attained iff all used routes
  have equal length, load and order count (\cref{lem:f4-range}).
\item \emph{Distance lower bound.}  Both of
  \[
    f_5 \;\ge\; \Bigl( \min_{(i,j)\in\Arcs} d_{ij} \Bigr)
                \sum_{i \in \Cust} z_i
    \qquad\text{and}\qquad
    f_5 \;\ge\; \frac{1}{2}\sum_{i\in\Cust} z_i
      \Bigl( \min_{j \ne i} d_{ij} + \min_{j \ne i} d_{ji} \Bigr)
  \]
  are valid, and neither uses the triangle inequality.
\end{enumerate}
\end{proposition}

\begin{proof}
(a) Sum \eqref{eq:c-cap} over $k$ and use \eqref{eq:c-assign}.  The rounded form
follows from integrality of $\sum_k u^k$.  (b) If at most $\kappa$ routes serve
$\sigma$ customers, the pigeonhole principle puts at least
$\lceil \sigma/\kappa \rceil$ of them on one route, whose diameter is by
\eqref{eq:Dr} at least the smallest diameter of a customer set of that
cardinality; the displayed bound then follows from \eqref{eq:c-Dmax}.  (c) is
\cref{lem:f4-range}.  (d) Each served customer has exactly one outgoing and one
incoming arc in $\widehat{G}$ by \eqref{eq:c-out}--\eqref{eq:c-in}, whose lengths
are at least the stated minima; each arc is counted at most twice.
\end{proof}

\begin{proposition}[Complexity]
\label{prop:hardness}
\textup{\textsc{MO-CVRP-OS}} is NP-hard, and deciding whether a feasible solution with
$f_1 = \sum_{i}o_i$, $f_2 \le \kappa$ and $f_5 \le B$ exists is NP-complete.
\end{proposition}

\begin{proof}
Restricting $\Feas$ to solutions with $\Ground(\Routes) = \Cust$ and
$|\Routes| \le \kappa$, and minimising $f_5$, is exactly the CVRP with $\kappa$
vehicles, which is NP-hard \cite{TothVigo2014}; membership in NP of the decision
version is clear since a solution is a polynomial-size certificate whose
objectives are computable in polynomial time.
\end{proof}

\subsection{Model size}
\label{sec:milp-size}

\begin{center}
\begin{tabular}{@{}lll@{}}
\toprule
Family & Count & Type \\
\midrule
$x^k_{ij}$   & $\overline{K}\,n(n+1)$ & binary \\
$y^k_i$      & $\overline{K}\,n$      & binary \\
$z_i$        & $n$                    & binary \\
$u^k$        & $\overline{K}$         & binary \\
$b_p$        & $P+1$                  & binary \\
$g_{ij}$     & $n(n+1)$               & continuous \\
$D^k, L^k, W^k, O^k$ & $4\overline{K}$ & continuous \\
$\Dmax,\Lmax,\Wmax,\Omax,\lambda$ & $5$ & continuous \\
$t^{L}_p, t^{W}_p, t^{O}_p$ & $3(P+1)$ & continuous \\
\midrule
\eqref{eq:c-Dk} (diameter)      & $\overline{K}\,\binom{n}{2}$ & \\
\eqref{eq:c-out}--\eqref{eq:c-in}, \eqref{eq:c-sym2} & $\Theta(\overline{K} n)$ & \\
\eqref{eq:c-flowcons}--\eqref{eq:c-flowreturn} & $\Theta(n^2)$ & \\
\eqref{eq:c-maxima}--\eqref{eq:c-maxima-ub}, \eqref{eq:c-balL-lin}--\eqref{eq:c-balO-lin}
  & $6\overline{K}+3$ & \\
\eqref{eq:c-mcL} and analogues  & $12(P+1)$ & \\
\bottomrule
\end{tabular}
\end{center}

With the default $\overline{K} = n$ the model has $\Theta(n^3)$ binary
variables and $\Theta(n^3)$ constraints, dominated by the arc variables and by
the diameter constraints \eqref{eq:c-Dk}.  The balance machinery, by contrast,
costs only $P+1$ additional binaries and $O(\overline{K} + P)$ constraints,
which is the payoff of the collapse \eqref{eq:f4-collapsed}.

This is a reference formulation.  Direct solution is realistic only for
instances with a few tens of customers, whereas the smallest instance in
\texttt{instances/} has $n = 100$ and the largest $n = \num{30000}$.  For those,
the formulation serves to define the problem unambiguously, to validate
heuristics on small sub-instances, and as the starting point for the
decompositions of \cref{sec:alt-formulations}.

% =========================================================================
\section{Modelling decisions and alternatives}
\label{sec:variants}

The informal specification of \textup{\textsc{MO-CVRP-OS}} leaves several points open.
This section states the decision taken in each case, and the alternatives that
were rejected, so that the model can be varied deliberately rather than by
accident.  \Cref{tab:decisions} summarises.

\begin{table}[htbp]
\centering
\small
\begin{tabular}{@{}>{\raggedright\arraybackslash}p{0.26\linewidth}>{\raggedright\arraybackslash}p{0.30\linewidth}>{\raggedright\arraybackslash}p{0.35\linewidth}@{}}
\toprule
Open point & Decision & Main alternatives \\
\midrule
Route diameter & customers only, \eqref{eq:Dr} &
  depot-inclusive; depot radius; longest leg (\cref{sec:alt-diameter}) \\
Balance combiner & egalitarian $\min$, \eqref{eq:beta} &
  weighted per-route composite; component-wise minima
  (\cref{sec:alt-balance}) \\
Balance normaliser & ratio to the maximum &
  range; standard deviation; Gini; Jain (\cref{sec:alt-balance}) \\
Length degeneracy & fifth objective $f_5$ &
  lexicographic tie-break; $\varepsilon$-constraint; sequence-optimal routes
  (\cref{sec:alt-f5}) \\
Service objective & $\sum_i o_i z_i$, $o_i = 1$ &
  delivered demand $\sum_i q_i z_i$; profits (\cref{sec:alt-f1}) \\
Connectivity & single-commodity flow &
  MTZ; generalised SECs; multi-commodity flow (\cref{sec:alt-sec}) \\
Route identification & vehicle indices &
  two-index; set partitioning (\cref{sec:alt-formulations}) \\
Fleet bound & $\overline{K} = n$ &
  any $\overline{K} < n$, at the cost of part of the front
  (\cref{sec:milp-data}) \\
$f_4$ exactness & grid of mesh $2^{-P}$ &
  fixed $\theta$; bisection; bilinear MIQCP (\cref{sec:alt-eps}) \\
\bottomrule
\end{tabular}
\caption{Modelling decisions.}
\label{tab:decisions}
\end{table}

\subsection{Route diameter}
\label{sec:alt-diameter}

\Cref{eq:Dr} measures the geographic spread of the \emph{customers} of a route.
Three alternatives are natural.

\begin{enumerate}[label=(\alph*),leftmargin=*]
\item \emph{Depot-inclusive diameter}
  $D'(r) = \max_{i,j \in C(r)\cup\{0\}} d_{ij}$, linearised by adding
  $D^k \ge d_{0i}\, y^k_i$ for $i \in \Cust$ to \eqref{eq:c-Dk}.  It never
  vanishes on a non-empty route (by \cref{as:depot}), so the degenerate
  ``one customer per route gives $f_3 = 0$'' behaviour of
  \cref{prop:bounds}(b) disappears and $f_3$ becomes
  $\max_r \max\{D(r), \max_{i\in C(r)} d_{0i}\}$.  This makes $f_3$ conflict
  with $f_2$ much more strongly, which may be desirable or not.
\item \emph{Depot radius} $\rho(r) = \max_{i \in C(r)} d_{0i}$, linearised by
  $\rho^k \ge d_{0i}\,y^k_i$.  It measures how far a route ventures rather than
  how spread out it is, and it is cheaper: $O(\overline{K} n)$ constraints
  instead of $O(\overline{K} n^2)$.
\item \emph{Longest leg} $\max_{(i,j) \in r} d_{ij}$, linearised by
  $D^k \ge d_{ij}\, x^k_{ij}$.  Unlike the other three, this depends on the
  visiting sequence, so \cref{prop:degeneracy} would not apply to it.
\end{enumerate}

The chosen definition \eqref{eq:Dr} has one property worth naming: it depends
on the \emph{partition} of the served customers only.  Minimising $f_3$ alone,
subject to serving everybody with at most $\kappa$ routes, is therefore the
min-max diameter clustering problem, which is NP-hard; $f_3$ contributes
genuine combinatorial difficulty on top of the routing.

\subsection{Combining the balance components}
\label{sec:alt-balance}

Write $\ell_r = L(r)/\Lmax$, $\upsilon_r = W(r)/\Wmax$ and
$\omega_r = O(r)/\Omax$ for the three normalised components of route $r$, and
let $\alpha,\beta,\gamma \ge 0$ with $\alpha+\beta+\gamma = 1$.  Three
combination rules were considered:
\begin{align}
  f_4^{\,\mathrm{egal}}(\Routes)
    &= \min_{r \in \Routes} \min\{\ell_r, \upsilon_r, \omega_r\}
    && \text{(adopted, \eqref{eq:beta})},
  \label{eq:bal-egal}\\
  f_4^{\,\mathrm{comp}}(\Routes)
    &= \min_{r \in \Routes}
       \bigl( \alpha \ell_r + \beta \upsilon_r + \gamma \omega_r \bigr)
    && \text{(per-route composite, then min)},
  \label{eq:bal-comp}\\
  f_4^{\,\mathrm{cw}}(\Routes)
    &= \alpha \min_{r} \ell_r + \beta \min_{r} \upsilon_r
       + \gamma \min_{r} \omega_r
    && \text{(component-wise minima, then combine)}.
  \label{eq:bal-cw}
\end{align}

\begin{proposition}
\label{prop:bal-order}
For every $\Routes \in \Feas$ and all admissible weights,
\begin{equation}
  f_4^{\,\mathrm{egal}}(\Routes)
  \;\le\; f_4^{\,\mathrm{cw}}(\Routes)
  \;\le\; f_4^{\,\mathrm{comp}}(\Routes),
\end{equation}
and all three equal $1$ simultaneously, precisely under the conditions of
\cref{lem:f4-range}.
\end{proposition}

\begin{proof}
For the right inequality, $\min_r$ of a sum is at least the sum of the
$\min_r$'s of its terms, applied to \eqref{eq:bal-comp} with non-negative
weights.  For the left, $\min\{a,b,c\} \le \alpha a + \beta b + \gamma c$
whenever $\alpha+\beta+\gamma=1$ and $\alpha,\beta,\gamma \ge 0$; apply this
with $a = \min_r \ell_r$ and so on, and note
$f_4^{\,\mathrm{egal}} = \min\{\min_r \ell_r, \min_r \upsilon_r,
\min_r \omega_r\}$ by \eqref{eq:f4-collapsed}.  All three equal $1$ iff each
ratio equals $1$ for every route, which is \cref{lem:f4-range}.
\end{proof}

The adopted rule \eqref{eq:bal-egal} is thus the most conservative of the
three, requires no weights to be chosen, and -- by the collapse
\eqref{eq:f4-collapsed} -- costs a single scalar variable $\lambda$.
\Cref{eq:bal-comp} is the most permissive: a route may be poor in one
dimension and compensate in another.  Implementing \eqref{eq:bal-comp} requires
one ratio variable $\pi_r$ per route and component, hence
$3\overline{K}(P+1)$ binaries instead of $P+1$; \eqref{eq:bal-cw} requires
three, which is still cheap.

Rules that are not ratios to the maximum are also possible: the range
$\Lmax - L^{\min}$ (to be minimised, and linear in the model variables, so it
would need no linearisation at all), the standard deviation of the route
lengths (convex quadratic), the Gini coefficient (linearisable with
$O(\overline{K}^2)$ absolute-value variables), or Jain's fairness index
$(\sum_k L^k)^2 / (\,|\Routes| \sum_k (L^k)^2)$ (a ratio of quadratics).  The
survey of \cite{Matl2018} compares such measures for vehicle routing
specifically and documents the pathology of \cref{prop:degeneracy} -- an equity
measure improved by making the best-off worse -- for min-max and range-based
measures alike, which is precisely why $f_5$ is retained here.

\subsection{The service objective}
\label{sec:alt-f1}

\Cref{eq:f1} counts orders.  With the default $o_i = 1$ of \cref{as:orders} it
counts served customers.  Two alternatives require no change to the
formulation, only to the coefficients of \eqref{eq:o1}:
\emph{delivered demand} $\sum_{i} q_i z_i$, appropriate when the quantity
shipped matters more than the number of stops, obtained by setting
$o_i = q_i$; and \emph{profit} $\sum_i p_i z_i$ for instance-specific prizes
$p_i$, which is the classical prize-collecting objective
\cite{Archetti2014}.  Note that the choice also propagates into the third
component of the balance measure \eqref{eq:beta}, which is defined in terms of
$O(r) = \sum_{i \in C(r)} o_i$; if orders and balance should be counted
differently, two separate parameter vectors must be introduced.

\subsection{Handling the length degeneracy}
\label{sec:alt-f5}

\Cref{prop:degeneracy} forces a choice.  Four responses are available.

\begin{enumerate}[label=(\alph*),leftmargin=*]
\item \emph{A fifth objective} (adopted).  The trade-off between balance and
  distance becomes explicit and is left to the decision maker.  Cost: a
  five-dimensional Pareto front, which is harder to compute and to present.
\item \emph{A lexicographic tie-break.}  Keep the four-objective vector and
  refine the dominance relation of \cref{def:dominance} by declaring
  $\Routes \prec \Routes'$ also when $F_{1:4}(\Routes) = F_{1:4}(\Routes')$ and
  $f_5(\Routes) < f_5(\Routes')$.  The resulting efficient set is a subset of
  the four-objective one containing, for each attainable
  $(f_1,f_2,f_3,f_4)$ vector, only the shortest solutions achieving it.  This
  is the least intrusive fix and is recommended when a four-objective front is
  required for comparability with other work.
\item \emph{An $\varepsilon$-constraint on distance.}  Keep four objectives and
  add $f_5 \le \overline{d}$ for an instance-specific budget, e.g.\ a multiple
  of a best-known CVRP solution value.  Sound, but the front then depends on an
  exogenous parameter.
\item \emph{Sequence-optimal routes.}  Redefine $\Feas$ to contain only
  solutions in which each route is a minimum-length Hamiltonian circuit on
  $C(r) \cup \{0\}$.  This removes the degeneracy at the source --
  \cref{prop:degeneracy} becomes vacuous because the sequence is a function of
  the customer set -- and reduces \textup{\textsc{MO-CVRP-OS}} to a pure partitioning
  problem.  The price is that feasibility becomes NP-hard to verify, so the
  condition has no compact MILP expression; it is natural only in the
  set-partitioning model of \cref{sec:alt-formulations}, where it amounts to
  restricting the column set $\Omega$ to sequence-optimal routes.
\end{enumerate}

Dropping $f_5$ from \cref{def:milp} recovers the pure four-objective problem;
nothing else in the formulation changes.

\subsection{Connectivity constraints}
\label{sec:alt-sec}

The single-commodity flow of \cref{sec:milp-flow} was chosen because it is
compact, because it is stronger than the Miller--Tucker--Zemlin alternative,
and because its validity argument survives optional service.  The alternatives:

\begin{enumerate}[label=(\alph*),leftmargin=*]
\item \emph{Miller--Tucker--Zemlin} \cite{MTZ1960}, per vehicle: variables
  $s^k_i$ with $q_i \le s^k_i \le Q$ and
  \[
    s^k_i - s^k_j + Q\,x^k_{ij} \;\le\; Q - q_j,
    \qquad i,j \in \Cust,\ i \ne j,\ k \in \Veh .
  \]
  Compact and simple, but its linear relaxation is dominated by that of
  \eqref{eq:c-flowcons}--\eqref{eq:c-flowcust}, and it needs
  $\Theta(\overline{K} n^2)$ constraints against $\Theta(n^2)$ for the
  aggregated flow.
\item \emph{Generalised subtour elimination constraints}, in the form adapted
  to optional service:
  \[
    \sum_{(i,j) \in \Arcs,\; i,j \in S} x_{ij}
      \;\le\; \sum_{i \in S} z_i - z_h,
    \qquad S \subseteq \Cust,\ |S| \ge 2,\ h \in S .
  \]
  Exponentially many, separated in polynomial time by minimum-cut
  computations.  Exact without \cref{as:demand}, which makes it the right
  choice if zero-demand customers must be supported
  (\cref{rem:demand-needed}).
\item \emph{Directed capacity constraints}
  $\sum_{(i,j) \in \delta^+(S)} x_{ij} \ge \frac{1}{Q}\sum_{i \in S} q_i z_i$
  for $S \subseteq \Cust$.  Valid as written; note that the familiar rounded
  form $\lceil q(S)/Q \rceil$ is \emph{not} directly available here, because
  the right-hand side depends on the variables $z$, and rounding is legitimate
  only after $z$ has been fixed (for example in a branch where the served set
  is decided).
\item \emph{Multi-commodity flow}, one commodity per served customer: the
  strongest compact relaxation, at $\Theta(n^3)$ variables, which is not
  attractive at the sizes considered here.
\end{enumerate}

\subsection{Alternative formulations}
\label{sec:alt-formulations}

\paragraph{Two-index model.}
Replacing $x^k_{ij}$ by $x_{ij}$ removes the factor $\overline{K}$ from the
model size and the vehicle symmetry with it, and suffices for $f_1$, $f_5$ and
connectivity.  It does not suffice for $f_2$, $f_3$ or $f_4$: the number of
routes is recoverable as $\sum_{j} x_{0j}$, but the diameter and the balance
measure need to know \emph{which} customers share a route, which the aggregated
arc vector does not express without reintroducing an assignment structure.  The
undirected variant, with $e_{ij} \in \{0,1\}$ for $i,j \in \Cust$ and
$e_{0j} \in \{0,1,2\}$ to allow single-customer routes, has the same
limitation.

\paragraph{Set partitioning.}
Let $\Omega$ be the set of all capacity-feasible routes, $a_{ir} = 1$ if route
$r \in \Omega$ visits customer $i$, and $\theta_r \in \{0,1\}$ select route $r$.
With $L_r$, $W_r$, $O_r$, $D_r$ the (constant) attributes of column $r$:
\begin{align}
  \sum_{r \in \Omega} a_{ir}\,\theta_r &= z_i \le 1, & &i \in \Cust, \\
  f_1 = \sum_{i} o_i z_i, \quad
  f_2 &= \sum_{r \in \Omega} \theta_r, \quad
  f_5 = \sum_{r \in \Omega} L_r \theta_r, \\
  \Dmax &\ge D_r\,\theta_r, & &r \in \Omega, \\
  \Lmax \ge L_r \theta_r, \quad
  \lambda \cdot \Lmax &\le L_r + \Lub\,(1-\theta_r), & &r \in \Omega,
\end{align}
and analogously for $W$ and $O$.  Connectivity and capacity are enforced
implicitly by the definition of $\Omega$, all four route attributes are
constants, and the sequence-optimality of \cref{sec:alt-f5}(d) is imposed
simply by restricting $\Omega$.  The obstacle is that $\Dmax$, $\Lmax$,
$\Wmax$, $\Omax$ and $\lambda$ \emph{couple} the columns, so the linear
relaxation is not a pure set-partitioning LP and its pricing problem is not the
usual elementary shortest path problem with resource constraints: the dual
information attached to the coupling constraints must be carried into the
pricing.  A practical route is to fix $(\Dmax, \Lmax, \Wmax, \Omax, \lambda)$
in an outer enumeration or bisection -- see \cref{sec:alt-eps} -- and to solve
a standard branch-and-price for each fixed tuple.

\subsection{Scalarisation, and exact treatment of \texorpdfstring{$f_4$}{f4}}
\label{sec:alt-eps}

Because $\Feas$ is combinatorial, the Pareto front has unsupported points and
weighted-sum scalarisation cannot enumerate it; the $\varepsilon$-constraint
method can \cite{Ehrgott2005}.  For \textup{\textsc{MO-CVRP-OS}} this is doubly
convenient.

\begin{enumerate}[label=(\alph*),leftmargin=*]
\item \emph{$f_4$ as a constraint is linear.}  Imposing $f_4 \ge \theta$ for a
  \emph{constant} $\theta \in [0,1]$ turns
  \eqref{eq:c-balL}--\eqref{eq:c-balO} into
  \[
    \theta\, \Lmax \le L^k + \Lub(1-u^k), \quad
    \theta\, \Wmax \le W^k + Q(1-u^k), \quad
    \theta\, \Omax \le O^k + \Oub(1-u^k),
  \]
  which are linear.  The entire apparatus
  \eqref{eq:c-lambda-exp}--\eqref{eq:c-balO-lin} can then be deleted, and the
  model is an exact MILP with no discretisation error whatsoever
  (\cref{rem:exactness}).  Since $f_4$ takes finitely many rational values
  (\cref{lem:integrality}), bisection on $\theta$ down to the separation
  $1/M^2$ of \cref{cor:resolution} computes $\max f_4$ exactly in
  $O(\log M)$ MILP solves, with $M = \max\{\Lub, Q, \Oub\}$.
\item \emph{$f_3$ as a constraint is a preprocessing rule.}  Imposing
  $f_3 \le \varepsilon_3$ is equivalent to forbidding two customers at distance
  more than $\varepsilon_3$ from sharing a route:
  \[
    y^k_i + y^k_j \le 1
    \qquad \text{for all } i<j \text{ with } d_{ij} > \varepsilon_3,\
    k \in \Veh,
  \]
  which also lets the variables $D^k$ and $\Dmax$ and the constraints
  \eqref{eq:c-Dk}--\eqref{eq:c-Dmax} be dropped entirely.  This is a far better
  behaved model than the min-max form and is the recommended way to explore the
  $f_3$ dimension.
\item \emph{Bilinear form.}  Keeping \eqref{eq:c-balL}--\eqref{eq:c-balO} as
  written yields a mixed-integer bilinear program (MIQCP) that a global solver
  can handle directly.  This is exact and compact but gives up the MILP
  machinery; it is mentioned for completeness, not recommended.
\item \emph{Dinkelbach.}  If only one of the three components of
  \eqref{eq:beta} is retained -- say the length component, so that
  $f_4 = L^{\min}/\Lmax$ -- then maximising $f_4$ is a single-ratio
  mixed-integer fractional program, solvable exactly by Dinkelbach's algorithm:
  iterate $\max\{L^{\min} - \theta\,\Lmax\}$ and update
  $\theta \leftarrow L^{\min}/\Lmax$ at the incumbent.  The sum-of-ratios
  structure of \eqref{eq:bal-comp} and \eqref{eq:bal-cw} does not admit this
  treatment, which is a further argument for the egalitarian rule
  \eqref{eq:bal-egal}.
\end{enumerate}

\subsection{Optional side constraints}
\label{sec:alt-side}

None of the following is part of \textup{\textsc{MO-CVRP-OS}} as defined, but each is a
one-line addition to \cref{def:milp} and each is standard in applications.

\begin{align}
  \sum_{i \in \Cust} o_i z_i &\;\ge\; \rho \sum_{i \in \Cust} o_i,
  & &\text{minimum service level } \rho \in [0,1]
  \text{ (see \cref{rem:trivial})}, \\
  z_i &\;=\; 1, \quad i \in M \subseteq \Cust,
  & &\text{mandatory customers}, \\
  \sum_{k \in \Veh} u^k &\;\le\; \kappa,
  & &\text{fleet cap}, \\
  L^k &\;\le\; \overline{L}^{\,\mathrm{route}}, \quad k \in \Veh,
  & &\text{maximum route length or duration}.
\end{align}
The first is the principled way to exclude the trivial efficient points of
\cref{rem:trivial}; note that it changes the ideal point and therefore any
normalisation built on it.

% =========================================================================
\bibliographystyle{plain}
\bibliography{references}

\end{document}
